% =====================================================================
% Secção 1 — Formulação do Problema
% Texto para o relatório (template IEEE). Baseado em docs/01_formulacao_problema/.
% Pacotes usados: nenhum além dos do template IEEEtran.
% Para acentuação correta, o preâmbulo deve ter:
%   \usepackage[utf8]{inputenc}  \usepackage[T1]{fontenc}  \usepackage[portuguese]{babel}
% =====================================================================

\section{Formulação do Problema}
\label{sec:formulacao}

Este trabalho utiliza o conjunto de dados \textit{Lending Club Loan Data (2007--2018)},
disponibilizado no Kaggle pela LendingClub, a maior plataforma norte-americana de
empréstimos \textit{peer-to-peer}. Os dados estão divididos em dois ficheiros:
\texttt{accepted\_2007\_to\_2018Q4.csv}, em que cada linha é um empréstimo financiado,
descrito por 151 colunas (perfil de crédito do mutuário, condições do empréstimo e
histórico de pagamentos), e \texttt{rejected\_2007\_to\_2018Q4.csv}, em que cada linha é
um pedido recusado, descrito por 9 colunas.

Todas as decisões desta secção --- tarefas, momentos de previsão, elegibilidade das
variáveis, coortes e divisão dos dados --- foram tomadas \textbf{antes} da análise
exploratória, com base apenas no enunciado e no dicionário de dados. O objetivo é evitar
que estas escolhas sejam influenciadas pelos resultados observados e, assim, prevenir a
fuga de informação (\textit{data leakage}).

% ---------------------------------------------------------------------
\subsection{Definição das tarefas}

Nas tarefas supervisionadas, a unidade de análise é um empréstimo financiado.

\textbf{Regressão.} Prever a taxa de juro anual (\texttt{int\_rate}) que a LendingClub
atribuiria a um pedido, no momento em que este é submetido e antes de a taxa ser fixada,
usando apenas as características do pedido e o perfil de crédito do requerente.
A população é o conjunto de todos os empréstimos financiados, uma vez que a taxa é
conhecida para todos desde a concessão.

\textbf{Classificação.} Prever, no momento da concessão do empréstimo --- quando as
condições (taxa, grau de risco e prestação) já estão fixadas, mas nenhum pagamento foi
observado --- o desfecho do empréstimo em três classes: \textit{Good standing},
\textit{At risk} e \textit{Default}.

\textbf{Clustering.} Identificar, sem prever qualquer alvo, grupos de mutuários com perfis
de crédito semelhantes e grupos de estados norte-americanos com comportamento semelhante
em termos de taxa de aprovação, montantes pedidos e \textit{risk scores}, combinando
pedidos aceites e rejeitados.

% ---------------------------------------------------------------------
\subsection{Momentos de previsão e informação disponível}

Uma variável só pode ser usada como preditor se o seu valor já for conhecido no momento
de previsão da tarefa. Para aplicar este critério, a vida de um empréstimo foi dividida
nos momentos da Tabela~\ref{tab:momentos}. A regressão é feita em T0 e a classificação
em T4; o clustering usa apenas o perfil do mutuário em T0.

\begin{table}[!t]
\caption{Momentos da vida de um empréstimo e informação que passa a existir}
\label{tab:momentos}
\centering
\footnotesize
\begin{tabular}{|p{0.22\columnwidth}|p{0.66\columnwidth}|}
\hline
\textbf{Momento} & \textbf{Exemplos de informação} \\
\hline
T0 Pedido & montante, prazo, finalidade, emprego, rendimento, DTI, localização, histórico de crédito (\textit{bureau}), campos do 2.º proponente \\
\hline
T1 Verificação & \texttt{verification\_status} \\
\hline
T2 \textit{Pricing} & \texttt{grade}, \texttt{sub\_grade}, \texttt{int\_rate}, \texttt{installment} \\
\hline
T3 Financiamento & \texttt{funded\_amnt}, \texttt{initial\_list\_status} \\
\hline
T4 Concessão & \texttt{issue\_d}, \texttt{disbursement\_method} \\
\hline
T5 Pagamentos & \texttt{total\_pymnt}, \texttt{recoveries}, \texttt{hardship\_*}, \texttt{settlement\_*}, \ldots \\
\hline
T6 Extração & \texttt{loan\_status}, \texttt{last\_fico\_range\_*}, \texttt{last\_credit\_pull\_d} \\
\hline
\end{tabular}
\end{table}

Destacam-se três consequências. Primeiro, \texttt{issue\_d} corresponde ao mês de
financiamento e não à data do pedido; por isso não é usado como preditor na regressão,
servindo apenas para a divisão temporal. Segundo, a antiguidade do histórico de crédito é
calculada relativamente a \texttt{issue\_d} na classificação, mas na regressão recorre-se a
variáveis de antiguidade já calculadas pelo \textit{bureau} (por exemplo,
\texttt{mo\_sin\_old\_rev\_tl\_op}), para não usar uma data futura. Terceiro, não sendo
possível garantir que a verificação do rendimento termina antes do \textit{pricing},
\texttt{verification\_status} é excluído da regressão. Assume-se que os campos do
\textit{bureau} refletem o relatório de crédito obtido no pedido, o que constitui uma
limitação deste instantâneo histórico.

% ---------------------------------------------------------------------
\subsection{Elegibilidade das variáveis}

Cada uma das 151 colunas foi classificada, para cada tarefa, como elegível, condicional ou
não elegível. A tabela completa encontra-se no repositório
(\texttt{docs/01\_formulacao\_problema/elegibilidade\_colunas.md}); a
Tabela~\ref{tab:elegibilidade} resume-a. Elegibilidade significa apenas que a utilização
da variável não constitui fuga de informação; a seleção final é feita mais tarde, com base
nos dados de treino.

\begin{table}[!t]
\caption{Resumo da elegibilidade das 151 colunas}
\label{tab:elegibilidade}
\centering
\footnotesize
\begin{tabular}{|l|c|c|c|c|}
\hline
\textbf{Tarefa} & \textbf{Elegível} & \textbf{Condicional} & \textbf{Excluída} & \textbf{Alvo} \\
\hline
Regressão & 92 & 4 & 54 & 1 \\
\hline
Classificação & 102 & 3 & 45 & 1 \\
\hline
Clustering & 74 & 1 & 76 & -- \\
\hline
\end{tabular}
\end{table}

As principais decisões foram:
\begin{itemize}
  \item \textbf{Desempenho pós-concessão} (38 colunas: pagamentos, recuperações,
  \texttt{hardship\_*}, \texttt{settlement\_*}, \texttt{last\_fico\_range\_*}): excluídas em
  todas as tarefas, bem como quaisquer variáveis delas derivadas.
  \item \textbf{Variáveis de \textit{pricing}} (\texttt{grade}, \texttt{sub\_grade},
  \texttt{installment}): excluídas da regressão, porque determinam ou permitem reconstruir
  a taxa; permitidas na classificação, por já estarem fixadas na concessão.
  \item \textbf{Histórico de crédito} (69 colunas): elegível nas três tarefas, incluindo
  campos como \texttt{acc\_now\_delinq} ou \texttt{chargeoff\_within\_12\_mths}, que
  descrevem outras contas do mutuário à data do pedido e não o empréstimo em análise.
  \item \textbf{Pedido conjunto} (\texttt{application\_type}, \texttt{*\_joint},
  \texttt{sec\_app\_*}): elegíveis; nos pedidos individuais, a ausência de valor é tratada
  como \textit{não aplicável} e não é substituída por zero.
  \item \textbf{Geografia} (\texttt{zip\_code}, \texttt{addr\_state}): não constitui fuga de
  informação, mas é condicionada a uma análise de \textit{fairness}, por poder funcionar
  como \textit{proxy} de características sensíveis.
  \item \textbf{Casos de processo}: na classificação, \texttt{funded\_amnt} e
  \texttt{initial\_list\_status} são permitidos, por serem conhecidos antes da concessão;
  \texttt{funded\_amnt\_inv} e \texttt{disbursement\_method} são excluídos, por o momento
  em que ficam definidos ser incerto.
\end{itemize}

% ---------------------------------------------------------------------
\subsection{Coorte e alvo}

O alvo da classificação é construído a partir de \texttt{loan\_status}, cujos valores
distintos foram listados (sem contagens) para definir o mapeamento da
Tabela~\ref{tab:alvo}. Os empréstimos ``\textit{Does not meet the credit policy}'',
concedidos ao abrigo de critérios de crédito anteriores, são incluídos e mapeados pelo seu
desfecho; esta indicação não é usada como preditor, por estar codificada na coluna do alvo.
Os empréstimos \textit{Current} são excluídos, por o seu desfecho ainda não ser conhecido.
A regressão não filtra por \texttt{loan\_status}.

\begin{table}[!t]
\caption{Construção do alvo da classificação}
\label{tab:alvo}
\centering
\footnotesize
\begin{tabular}{|p{0.25\columnwidth}|p{0.63\columnwidth}|}
\hline
\textbf{Classe} & \textbf{Valores de \texttt{loan\_status}} \\
\hline
Good standing & Fully Paid; Does not meet the credit policy. Status:Fully Paid \\
\hline
At risk & In Grace Period; Late (16--30 days); Late (31--120 days) \\
\hline
Default & Charged Off; Default; Does not meet the credit policy. Status:Charged Off \\
\hline
Excluído & Current; sem valor \\
\hline
\end{tabular}
\end{table}

A exclusão dos \textit{Current} introduz um viés de seguimento desigual
(\textit{unequal follow-up}): nas coortes recentes, os empréstimos já terminados são
sobretudo os que entraram em incumprimento cedo ou foram pagos antecipadamente. Este
problema é resolvido no desenho da divisão temporal.

% ---------------------------------------------------------------------
\subsection{Desenho da divisão temporal}

As partições usam cortes de calendário explícitos sobre \texttt{issue\_d}, no início de cada
mês, de modo que empréstimos do mesmo mês ficam sempre na mesma partição. Define-se um
\textit{cutoff} de observação comum a 31/12/2018, justificado pela cobertura documentada do
conjunto de dados (até ao 4.º trimestre de 2018) e não inferido a partir das datas de
pagamento. Simula-se um modelo treinado a 01/01/2015 e aplicado aos empréstimos concedidos
em 2015. A Tabela~\ref{tab:split} resume as partições.

\begin{table}[!t]
\caption{Partições temporais por tarefa (\texttt{issue\_d})}
\label{tab:split}
\centering
\footnotesize
\begin{tabular}{|p{0.22\columnwidth}|p{0.30\columnwidth}|p{0.30\columnwidth}|}
\hline
\textbf{Partição} & \textbf{Regressão} & \textbf{Classificação} \\
\hline
Treino e validação & 2007-06 a 2014-12 & 2007-06 a 2011-12 \\
\hline
Embargo & -- & 2012-01 a 2014-12 \\
\hline
Teste final & 2015 & 2015 \\
\hline
Não usado & 2016--2018 & 2016--2018 \\
\hline
Coorte & Todos os financiados & 36 meses, não \textit{Current}, prazo terminado \\
\hline
\end{tabular}
\end{table}

Na regressão, o alvo é conhecido na concessão, pelo que todos os rótulos de treino eram
observáveis na data simulada. A validação é feita por validação cruzada temporal com
janela expansível por ano civil.

Na classificação, um rótulo só é usado se o empréstimo já tiver atingido o fim do prazo
(\texttt{issue\_d} + \texttt{term}) na data relevante. O treino inclui, assim, os
empréstimos de 36 meses concedidos até 2011-12, cujo prazo terminou antes da data de treino
simulada; os concedidos entre 2012 e 2014 ficam num período de embargo, por ainda não terem
terminado em 01/01/2015; e o teste final usa os concedidos em 2015, que terminam antes do
\textit{cutoff}. Esta regra garante que os rótulos de treino eram observáveis na data
simulada e que todos os empréstimos têm um seguimento completo, eliminando o viés referido.
Como os empréstimos de 60 meses de 2015 só terminariam em 2020, a classificação fica
definida apenas para empréstimos de 36 meses.

A análise exploratória, os testes de hipóteses e o clustering usam apenas o período de
desenvolvimento (2007--2014); as análises que envolvem o desfecho usam apenas o treino da
classificação. O teste final é usado uma única vez, no fim do trabalho. Todas as
transformações aprendidas a partir dos dados são ajustadas apenas no treino de cada
\textit{fold}; eventuais reamostragens são aplicadas apenas dentro dos \textit{folds} de
treino, e as amostras são aleatórias, com semente fixa.

% ---------------------------------------------------------------------
\subsection{Limitações}

\begin{itemize}
  \item Os campos do pedido e do \textit{bureau} são assumidos como registados no momento do
  pedido; \texttt{loan\_amnt} pode refletir reduções feitas pelo departamento de crédito.
  \item A classe \textit{At risk} corresponde ao estado do empréstimo na data de extração e
  não a um desfecho final; com a regra do prazo terminado, torna-se rara.
  \item O rótulo usado é o estado a 31/12/2018 e não o estado exato na data de treino
  simulada.
  \item A classificação usa um treino reduzido (2007--2011) e é avaliada num período
  posterior, depois de mudanças na política de crédito e no mercado.
  \item A identificação de mutuários repetidos depende da disponibilidade de
  \texttt{member\_id}, verificada na limpeza.
\end{itemize}
